% ---------------------------------------------------------------------
\section{重夸克质量与变味方案}\label{sec:disall-heavy}

粲与底夸克的质量 $m$ 介于 $\Lambda$ 与 $Q$ 之间：$Q\sim m$ 时重夸克只在末态成对产生，$Q\gg m$ 时 $\as\ln(Q^2/m^2)$ 变大，重夸克又应当作为部分子参与演化。
$\MSbar$ 与质量无关，重夸克不会自动退耦，退耦要通过 $n_f$ 与 $n_f-1$ 味有效理论的匹配来实现（\cref{rem:qcd-decoupling}）。下面 $Q$ 同时表示重夸克（如\cref{fig:disall-heavy}）与 $Q^2=-q^2$，从上下文可以区分。

\begin{theorem}[含重夸克质量的因子化（Collins 1998） \tagA]\label{thm:disall-heavy}
在硬部分中保留重夸克质量 $m_h$，并按 CWZ/ACOT 框架引入重夸克分布\cite{CollinsWilczekZee1978,AivazisEtAl1994}，DIS 结构函数仍取因子化形式
\begin{equation}\label{eq:disall-heavy}
\frac{F_a}{x}=\sum_i f_i\otimes\hat C_{a,i}\Big(\frac{m_h}Q,\frac Q\mu,\as(\mu)\Big)+R_a,
\end{equation}
其中 $i$ 包括重夸克，余项 $R_a$ 被 $\Lambda/Q$ 的某个幂次一致地压低，与 $m_h/Q$ 的大小无关：\cite{Collins1998b}的摘要写道
「“The power-suppressed corrections to the factorization formula are uniformly suppressed by a power of $\Lambda/Q$, independently of the size of heavy quark masses, $M$, relative to $Q$.”」
非微扰地位：\cref{asm:intro-bridge}（MA1、MA2）。
\uses{thm:disall-factorization,rem:qcd-decoupling,asm:intro-bridge}
% errtype: E-S
\end{theorem}

\begin{proofsketch}
论证与\cref{thm:disall-factorization}相同，只是重夸克质量保留在硬部分与相应的近似子中：$m_h\sim Q$ 时含重夸克的线是硬的；$m_h\ll Q$ 时重夸克也可以是活动部分子，其共线发散被 $m_h$ 截断为 $\ln(Q^2/m_h^2)$。
重整化采用 CWZ 的退耦减除方案\cite{CollinsWilczekZee1978}：$\mu<m_h$ 时理论等价于 $n_f-1$ 味的有效理论，$\mu>m_h$ 时用 $n_f$ 味的 $\MSbar$（\cref{rem:qcd-decoupling}）。
所得系数在 $Q\sim m_h$ 时趋于固定味方案（FFNS），在 $Q\gg m_h$ 时趋于零质量变味方案；一般质量变味方案（ACOT\cite{Aivazis:1993kh,AivazisEtAl1994}、S-ACOT\cite{Kramer:2000hn}、TR\cite{Thorne:1997ga}、
BMSN\cite{Buza:1996wv}、FONLL\cite{Forte:2010ta}）减去两者的重叠，在所计算的阶上是等价的处方（\cref{ex:disall-heavy-lo}）。完整论证见\cite{Collins1998b}。
\end{proofsketch}

\paragraph{地位与缺口}
本定理继承\cref{thm:disall-factorization}的全部瓶颈（\cref{asm:appN-collins,asm:appN-ward-product,asm:appN-cutsum}等）与具名缺口\cref{asm:appN-brst-nonlocal}；质量的保留不引入新的未证假设，余项的一致性同样是标度陈述（E-S）。
这一结论是对只有一个强子的 DIS 而言的：对强子—强子碰撞，\cite[第~14.6~节，第~569--570~页]{Collins2011}指出因子化只保证领头幂次上红外发散的相消，形如 $(m^2/Q^2)\times(1/\varepsilon)$ 的项可能残留，
在硬散射中保留入射重夸克的质量是“a bad idea”；把 ACOT 一类的方案移植到强子—强子碰撞时须记住这一区别。

\begin{example}[领头阶 $\gamma^*g\to Q\bar Q$]\label{ex:disall-heavy-lo}
考虑 $\gamma^*(q)+g(k)\to Q(p_1)+\bar Q(p_2)$，$p_1^2=p_2^2=m^2$（\cref{fig:disall-heavy}(a)），$z\equiv Q^2/(2k\cdot q)$。记 $\epsilon_m\equiv m^2/Q^2$，$\hat s=(k+q)^2=Q^2(1-z)/z$，$\beta\equiv\sqrt{1-4m^2/\hat s}=\sqrt{1-4\epsilon_mz/(1-z)}$，阈值为 $z\le1/(1+4\epsilon_m)$。
与无质量情形相比，旋量求和改为 $\sla p_1+m$、$\sla p_2-m$，传播子分母为 $t_1\equiv(k-p_1)^2-m^2=-\frac{\hat s+Q^2}2(1-\beta\cos\theta)$ 与 $u_1\equiv(k-p_2)^2-m^2=-\frac{\hat s+Q^2}2(1+\beta\cos\theta)$（$\gamma^*g$ 质心系，$\theta$ 为 $\vec p_1$ 与 $\vec k$ 的夹角），
极化求和仍可取 $-g_{\rho\sigma}$（我们以物理极化求和核对过），四维两体相空间为 $\frac\beta{16\pi}\int_{-1}^1\dd\cos\theta$。用\cref{def:disall-breit}的投影（$D=4$）并对 $\cos\theta$ 积分，得到（本书的推导）
\begin{equation}\label{eq:disall-heavy-lo}
\begin{aligned}
\frac{\hat F^g_2}{ze_Q^2}&=\frac{\as}{2\pi}\,2T_F\Big\{\big[z^2+(1-z)^2+4\epsilon_mz(1-3z)-8\epsilon_m^2z^2\big]\ln\frac{1+\beta}{1-\beta}\\
&\hspace{8.2em}+\beta\big[8z(1-z)-1-4\epsilon_mz(1-z)\big]\Big\},\\
\frac{\hat F^g_L}{ze_Q^2}&=\frac{\as}{2\pi}\,2T_F\Big\{-8\epsilon_mz^2\ln\frac{1+\beta}{1-\beta}+4\beta z(1-z)\Big\}
\end{aligned}
\end{equation}
（阈值以下为零）。$F_2$ 的结果与\cite[式~(49)--(51)]{Forte:2010ta}一致（那里 $a_s=\as/2\pi$，因子 2 同样计入 $Q$ 与 $\bar Q$；无质量极限见该文式~(52)--(53)），$F_L$ 的结果以同一方法得到。对数 $\ln\frac{1+\beta}{1-\beta}$ 来自 $1/t_1$、$1/u_1$ 在 $\cos\theta\to\pm1$ 附近的增强：$m$ 正规化了重夸克与胶子共线的区域。
当 $m\to0$ 时 $\ln\frac{1+\beta}{1-\beta}=\ln\frac{Q^2}{m^2}+\ln\frac{1-z}z+\order{\epsilon_m}$，于是
$\hat F^g_2/(ze_Q^2)\to\frac{\as}{2\pi}2[P_{qg}\ln\frac{Q^2}{m^2}+C_2^g]$，$\hat F^g_L/(ze_Q^2)\to\frac{\as}{2\pi}2C_L^g$，修正为 $\order{\epsilon_m\ln\epsilon_m}$：
无质量计算中的极点 $-\frac2\varepsilon P_{qg}$ 被 $2P_{qg}\ln(Q^2/m^2)$ 代替，有限部分恰为 $\MSbar$ 的 $C_2^g$（\cref{prop:disall-nlo}）。另一方面，由\cref{def:pdf-quark}在光锥规范中直接计算单圈的重夸克—胶子分布
（$\bar Q$ 在壳、$Q$ 离壳的割图，$D$ 维横动量积分，$\MSbar$ 扣除紫外极点）得 $f_{Q/g}(z,\mu)=\frac{\as}{2\pi}P_{qg}(z)\ln\frac{\mu^2}{m^2}+\order{\as^2}$，没有附加常数。
于是\cref{fig:disall-heavy}(b) 的减除项（Born 取无质量形式，即 S-ACOT 的选择\cite{Kramer:2000hn}；$\mu<m$ 时按 CWZ 方案 $f_{Q/g}=0$，不作减除）为 $[f_{Q/g}+f_{\bar Q/g}]\otimes\hat F^{(0)}_2/z=e_Q^2\frac{\as}{2\pi}2P_{qg}\ln(\mu^2/m^2)$，减除后的胶子系数
$\hat C_{2,g}=\frac{\as}{2\pi}[2T_F\{\cdots\}-2P_{qg}\ln\frac{\mu^2}{m^2}]$（$\{\cdots\}$ 为\cref{eq:disall-heavy-lo}第一式的花括号）在 $\mu=Q$、$m\to0$ 时趋于 $\frac{\as}{2\pi}2C_2^g$；它与光滑函数的卷积对一切 $m/Q$ 一致有界（阈值以上未被抵消的 $\ln(Q^2/m^2)$ 只占宽度 $\order{\epsilon_m}$ 的区间）。
数值上，$z=0.1$ 时 $\{\cdots\}$ 与其无质量渐近式 $C_2^g/T_F+(P_{qg}/T_F)\ln(Q^2/m^2)$ 之差在 $\epsilon_m=0.1,\,0.01,\,0.001$ 时依次为 $0.074,\,0.014,\,0.002$。
在 $\order{\as}$，重夸克初态项 $e_Q^2(f_Q+f_{\bar Q})$ 加上减除后的胶子项，恰好回到固定味方案的结果：各方案在所计算的阶上等价，变味方案的好处是 $f_Q$ 的演化对 $\as\ln(\mu^2/m^2)$ 作了重求和。
\uses{thm:disall-heavy,def:disall-breit,prop:disall-nlo,def:pdf-quark}
\end{example}

%%% BEGIN fig:disall-heavy
\begin{figure}[htbp]
\centering
\begin{tikzpicture}[fd,
  fb lab/.style={font=\footnotesize,inner sep=1.5pt},
  fb heavy/.style={draw,solid,line width=1.5pt,postaction={decorate,solid,decoration={markings,
      mark=at position #1 with {\arrow[xshift=3.6pt,line width=0.8pt]{Stealth[length=7pt,width=6pt]}}}}},
  fb heavy/.default=0.5]
  % ======================= (a) LO photon--gluon fusion with massive quarks =======================
  \begin{scope}
    \coordinate (B) at (0,0.75); \coordinate (A) at (0,-0.75);
    \draw[photon] (-1.6,0.75) -- (B);
    \draw[gluon] (-1.6,-0.75) -- (A);
    \draw[fb heavy] (B) -- (1.55,0.75);
    \draw[fb heavy] (A) -- (B);
    \draw[fb heavy] (1.55,-0.75) -- (A);
    \node[vertex] at (A) {}; \node[vertex] at (B) {};
    \path[mom={$q$},mom sep=0.26cm] (-1.6,0.75) -- (B);
    \path[mom'={$k$},mom sep=0.3cm] (-1.6,-0.75) -- (A);
    \path[mom={$p_1$},mom sep=0.28cm] (B) -- (1.55,0.75);
    \path[mom'={$p_2$},mom sep=0.3cm] (A) -- (1.55,-0.75);
    \node[fb lab,anchor=east] at (-1.62,0.75) {$\gamma^*$};
    \node[fb lab,anchor=east] at (-1.62,-0.75) {$g$};
    \node[fb lab,anchor=west] at (1.57,0.75) {$Q$};
    \node[fb lab,anchor=west] at (1.57,-0.75) {$\bar Q$};
    \node[fb lab,anchor=west,align=left] at (2.3,0) {$+$\ 交叉图\\\phantom{$+$\ }（$p_1\leftrightarrow p_2$）};
    \node[font=\small,anchor=north] at (0.85,-1.45) {(a)\ 领头阶 $\gamma^*g\to Q\bar Q$（保留 $m$）};
  \end{scope}
  % ======================= (b) ACOT/CWZ subtraction term =======================
  \begin{scope}[xshift=8.15cm]
    \coordinate (B) at (0,0.75);
    \node[jet blob,ellipse,minimum width=1.55cm,minimum height=0.9cm,inner sep=0pt] (F) at (0,-0.75) {};
    \draw[gluon] (-1.75,-0.75) -- (F.west);
    \draw[photon] (-1.6,0.75) -- (B);
    \draw[fb heavy=0.55] (F.north) -- (B);
    \draw[fb heavy] (B) -- (1.55,0.75);
    \draw[fb heavy=0.6] (1.6,-1.25) -- (F.-25);
    \node[vertex] at (B) {};
    \node[font=\footnotesize] at (F) {$f_{Q/g}$};
    \path[mom={$q$},mom sep=0.26cm] (-1.6,0.75) -- (B);
    \path[mom'={$k$},mom sep=0.3cm] (-1.75,-0.75) -- (F.west);
    \node[fb lab,anchor=east] at (-1.62,0.75) {$\gamma^*$};
    \node[fb lab,anchor=east] at (-1.77,-0.75) {$g$};
    \node[fb lab,anchor=west] at (1.57,0.75) {$Q$};
    \node[fb lab,anchor=west] at (1.62,-1.25) {$\bar Q$};
    \node[fb lab,anchor=east] at (-0.06,0.02) {$Q$};
    \node[fb lab,anchor=west,text=gray!30!black] at (0.14,0.38) {$\hat\sigma(\gamma^*Q\to Q)$};
    \node[font=\small,anchor=north] at (0,-1.45) {(b)\ 减除项 $f_{Q/g}\otimes\hat\sigma(\gamma^*Q\to Q)$};
  \end{scope}
  \node[font=\Large] at (5.45,0) {$-$};
\end{tikzpicture}
\caption{含重夸克质量的 DIS（粗线为质量 $m$ 的重夸克 $Q$）。(a)~领头阶的光子—胶子融合，保留完整的质量依赖，两图之和的平方给出领头阶的胶子道系数；
(b)~变味（ACOT/CWZ）方案中的减除项 ${f_{Q/g}\otimes\hat\sigma(\gamma^*Q\to Q)}$：胶子先在重夸克分布 $f_{Q/g}$（蓝色，靶共线）中分裂为 $Q\bar Q$，其中的 $Q$ 再发生领头阶散射 $\gamma^*Q\to Q$。
二者之差在 $m\to0$ 时抵消 (a) 中的共线对数 $\ln(Q^2/m^2)P_{qg}$，给出对 $m/Q$ 一致有效的系数函数（\cref{thm:disall-heavy}，\cref{ex:disall-heavy-lo}）。}
\label{fig:disall-heavy}
\end{figure}
%%% END fig:disall-heavy
